\subsection{Ebene durch drei Punkte}
\label{subsec:analytic_geometry_plane_through_three_points}

Drei nicht auf einer Geraden liegende Punkte bestimmen eine Ebene eindeutig.
Verbindungsvektoren zwischen den Punkten liefern unmittelbar zwei geeignete
Richtungen.

\begin{TheoremBox}{Ebene durch drei Punkte}{analytic_geometry_plane_from_three_points}
  Sind $A$, $B$ und $C$ nicht \emph{kollinear} \EnglishTerm{collinear}, so sind beispielsweise
  $\overrightarrow{AB}$ und $\overrightarrow{BC}$ linear unabhängig. Damit
  gilt
  \[
    E:\quad
    \vec x=\vec a
    +\lambda\overrightarrow{AB}
    +\mu\overrightarrow{BC},
    \qquad \lambda,\mu\in\R,
  \]
  wobei $\vec a=\overrightarrow{OA}$ der Ortsvektor von $A$ ist.
\end{TheoremBox}

\begin{center}
  \CachedTikzFigure{content/resources/chapters/analytic_geometry/tikz/plane_through_three_points.tex}
\end{center}

\begin{ExampleBox}{Ebene durch drei Punkte bestimmen}{analytic_geometry_plane_from_points_example}
  Gesucht ist die Ebene durch
  \[
    A(6,3,-5),\qquad B(2,-1,3),\qquad C(1,2,4).
  \]
  Wir wählen $A$ als Aufpunkt und konstruieren zwei Richtungen:
  \begin{ClozeArea}[34]{10}
    \[
      \begin{aligned}
        \overrightarrow{AB}
          &=\begin{pmatrix}2\\-1\\3\end{pmatrix}
            -\begin{pmatrix}6\\3\\-5\end{pmatrix}
            =\begin{pmatrix}-4\\-4\\8\end{pmatrix}
            =-4\begin{pmatrix}1\\1\\-2\end{pmatrix},\\[1em]
        \overrightarrow{BC}
          &=\begin{pmatrix}1\\2\\4\end{pmatrix}
            -\begin{pmatrix}2\\-1\\3\end{pmatrix}
            =\begin{pmatrix}-1\\3\\1\end{pmatrix}.
      \end{aligned}
    \]
  \end{ClozeArea}
  Die beiden Richtungen sind keine Vielfachen voneinander. Eine kompakte
  Parameterdarstellung lautet daher
  \[
    E:\quad
    \vec x
    =\begin{pmatrix}6\\3\\-5\end{pmatrix}
    +\lambda\begin{pmatrix}1\\1\\-2\end{pmatrix}
    +\mu\begin{pmatrix}-1\\3\\1\end{pmatrix},
    \qquad \lambda,\mu\in\R.
  \]
\end{ExampleBox}
